% ============================================================================
%  第 6 章 实验 —— body/06-experiments.tex
%  职责：可复现。设置 → 基线 → 结果（图/表）→ 消融 → 效率。
% ============================================================================

\section{实验与结果}
\label{sec:experiments}

% ---------------------------------------------------------------------------
\subsection{实验设置}

\begin{table}[t]
  \centering
  \caption{主要实验参数}
  \label{tab:settings}
  \begin{threeparttable}
  \begin{tabular}{@{}llcc@{}}
    \toprule
    类别 & 参数 & 取值 & 单位 \\
    \midrule
    空域   & 栅格分辨率     & \num{50}            & \si{\metre} \\
    时间   & 时间步长       & \num{60}            & \si{\second} \\
    风险   & 权重 $w_1$/$w_2$/$w_3$ & 0.5/0.2/0.3 & -- \\
    算法   & 标签上限 $q$   & \num{3}             & -- \\
    \bottomrule
  \end{tabular}
  \begin{tablenotes}
    \footnotesize \item[] 注：权重组合的敏感性分析见 \refsec{sec:sensitivity}。
  \end{tablenotes}
  \end{threeparttable}
\end{table}

基线方法选取 \basestatic（静态风险场）与 Greedy-Risk（贪心）。

% ---------------------------------------------------------------------------
\subsection{合成场景}

\begin{figure}[t]
  \centering
  \begin{subfigure}[b]{0.48\linewidth}
    \centering
    \placeholderfig{静态风险场下的规划结果}
    \caption{静态风险场}
    \label{fig:synthetic-static}
  \end{subfigure}
  \hfill
  \begin{subfigure}[b]{0.48\linewidth}
    \centering
    \placeholderfig{本文方法结果}
    \caption{\methodnamelatex（本文）}
    \label{fig:synthetic-ours}
  \end{subfigure}
  \caption{合成场景中不同时刻的规划路径对比。}
  \label{fig:synthetic}
\end{figure}

如 \reffig{fig:synthetic} 所示，……\todo{补 2～3 句结果分析，务必给出定量指标。}

% ---------------------------------------------------------------------------
\subsection{真实城市场景}

% ============================================================================
%  主结果表 —— assets/tables/main-results.tex
%  位置：第 \ref{sec:experiments} 节
% ============================================================================

\begin{table}[t]
  \centering
  \caption{真实城市场景下的主要结果（均值 $\pm$ 标准差；风险类指标越低越好，路径长度自然越低越好）}
  \label{tab:main-results}
  \begin{threeparttable}
  \begin{tabular}{@{}lcccc@{}}
    \toprule
    方法 & 综合风险 & 致死风险 & 噪声影响 & 路径长度 (km) \\
    \midrule
    \basestatic     & xx.xx $\pm$ x.xx & xx.xx & xx.xx & xx.xx \\
    Greedy-Risk     & xx.xx $\pm$ x.xx & xx.xx & xx.xx & xx.xx \\
    \textbf{本文}   & \textbf{xx.xx} $\pm$ x.xx & \textbf{xx.xx} & xx.xx & xx.xx \\
    \midrule
    相对最优基线   & $-$\,xx.x\% & $-$\,xx.x\% & $-$\,xx.x\% & $+\,$x.x\% \\
    \bottomrule
  \end{tabular}
  \begin{tablenotes}
    \footnotesize
    \item[] 注：最优值加粗显示；相对值为各方法的均值对比。
  \end{tablenotes}
  \end{threeparttable}
\end{table}


% ---------------------------------------------------------------------------
\subsection{消融实验}

依次移除时间维传播、累积危险率形式与支配剪枝，观察各项指标变化，
验证各模块贡献。

% ---------------------------------------------------------------------------
\subsection{计算效率}

记录各方法在相同地图上的运行时间与扩展节点数：在不显著增加运行时间的
前提下，本文方法将……$X\%$。

% ---------------------------------------------------------------------------
\subsection{敏感性分析}
\label{sec:sensitivity}

考察权重 $w_i$、栅格分辨率与时间步长对结果的影响。
